\begin{tikzpicture}
    % === DEFINICIÓN DE ESPACIO ===
    \coordinate (A) at (5.715,5.045);
    % (Líneas de ayuda para cuadricular el espacio durante el diseño)
    % \draw[help lines,xstep=0.1,ystep=0.1] (0,0) grid (10,6);
    % \draw[help lines,xstep=1,ystep=1,black] (0,0) grid (10,6);
    
    % === INCLUSIÓN DE ILUSTRACIONES TÉCNICAS EXTERNAS ===
    % Base del vehículo grúa anclada a la izquierda
    \node[anchor=south west, inner sep=0] (image) at (0,0) {\includegraphics[width=6cm]{grua.pdf}};
    % Silueta del operador anclada a la derecha
    \node[anchor=south west, inner sep=0] (image) at (8.65,0) {\includegraphics[width=4.55mm]{persona.pdf}};
    
    % === TRAZADO DEL SISTEMA DE TENSIONES (CABLES) ===
    % 1. Eje vertical de referencia desde la polea A
    \draw[dashed, thin, gray] (A) -- ++(0,-2) coordinate (num);

    % 2. Cable principal (polea a bloque)
    \draw[teal!80, line width=0.3mm] (A) -- ++(-85:2.65) coordinate (C);

    % 3. Eje horizontal de referencia en el nudo C
    \draw[dashed, thin, gray] (C) -- ++(2,0) coordinate (al);

    % 4. Enganche inferior hacia la masa
    \draw[teal!80, line width=0.3mm] (C) -- ++(0,-0.8) coordinate (D);

    % 5. Cable de tracción lateral (hacia el operador)
    \draw[teal!80, line width=0.3mm] (C) -- ++(-34.5:3.4) coordinate (B);

    % === NODOS Y ETIQUETAS VECTORIALES ===
    \fill[teal!80,nodes={black, font=\footnotesize}] (A) node[right] {$A$} circle (1pt);
    \fill[teal!80,nodes={black, font=\footnotesize}] (C) node[left] {$C$} circle (1pt);
    \fill[teal!80,nodes={black, font=\footnotesize}] (B) node[below left] {$B$} circle (1pt);

    % === CARGA Y MASA ESTÁTICA ===
    % Detalle del conector/enganche
    \draw[draw=black, rounded corners=0.5mm, fill=white] ($(D) - (0.1,0)$) -- ($(D) - (0.1, 0.2)$) -- ($(D) + (0.1, -0.2)$) -- ($(D) + (0.1, 0)$) -- cycle;

    % Bloque trapezoidal (Masa de 1200 lb)
    \filldraw[draw=black, fill=colorFondo!40, rounded corners=1mm] ($(D)-(0.3,0.2)$) -- ++(0.6,0) -- ++(0.2,-0.7) -- ++(-1,0) -- cycle;

    % Texto de magnitud
    \node[font=\tiny] at ($(D)+(0,-0.55)$) {\qty{1200}{\lb}};

    % === COTAS DE ÁNGULO E INCLINACIÓN ===
    % Desviación del cable principal respecto a la vertical
    \pic [draw=colorCotas, nodes={anchor=east}, "\ang{5}", angle radius=1.7cm, angle eccentricity=1, font=\footnotesize\color{colorCotas}] {angle = num--A--C};

    % Ángulo de tiro del operador respecto a la horizontal
    \pic [draw=colorCotas, latex-, "$\alpha$", angle radius=1cm, angle eccentricity=1.2, font=\small\color{colorCotas},] {angle = B--C--al};
\end{tikzpicture}